\section{背景と関連研究}

\subsection{バイナリレベルのTPL検出}
バイナリレベルのTPL検出は，対象バイナリから文字列，命令，関数の意味的特徴，関数呼出し関係などを抽出し，既知ライブラリとの対応を求める．ModX\cite{modx}は意味的特徴とプログラムのモジュール化を組み合わせて部分的に取り込まれたTPLを検出し，BinaryAI\cite{binaryai}は関数の意味的表現とリンク時の局所性を用いて候補を絞り込む．これらはコンパイラや最適化による表現の違いを越えて候補を探すことを目的としており，候補となる関数対をどれだけ高精度に見つけ，ライブラリ単位へ適切に集約するかに主眼を置く．一方，候補となった関数対が実際に同じ入力に対して同じ出力を返すかどうかは，特徴や構造の類似度からは直接確認できない．本研究は，これらの手法が出力する候補を対象に，実行結果の一致を後段で再検証する．

\subsection{バイナリの意味比較}
BinHunt\cite{binhunt}は基本ブロックをシンボリック実行し，意味的に対応するブロックと制御フローから関数の類似性を評価する手法であり，iBinHunt\cite{ibinhunt}はこれを関数間制御フローの差分探索へ拡張した．これらはバイナリ全体または関数全体を対象に意味的な差分を探索するものであり，既に候補として与えられた関数対を検証する後段処理としては設計されていない．

BinGo\cite{bingo}はクロスアーキテクチャでの類似関数検索を扱い，BLEX\cite{blex}とBinSim\cite{binsim}は実行またはトレースに基づく意味比較を行う．これらは動的解析を前提とするため，実行環境の準備や網羅的な入力生成が必要になる．本研究は静的解析であるHigh P-codeとSMTソルバを用いることで，実行環境なしに入出力の対応関係を検証する．

DDEC\cite{ddec}は，x86の2つのループが同じ動作をするかを検証する手法であり，実行結果から対応関係の候補を求め，SMTで検証する．低水準コード等価性検証\cite{lowlevel-eq}は，低水準コードの意味的一致を扱う．これらはループを含むプログラムの等価性判定という点で本研究と関係するが，対象は個別のループやプログラム断片であり，TPL検出の後段処理として大量の関数対応候補を検証する枠組みではない．

High P-codeの形式意味論は既存研究で検討されている\cite{formal-pcode}．これはHigh P-codeの意味を定義するものであり，関数間の等価性判定手法ではないが，本研究がHigh P-codeを解析対象とする際の意味論的な裏付けとなる．

\subsection{既存研究の課題と本研究の位置づけ}
バイナリレベルのTPL検出研究は，候補生成とライブラリ単位への集約の精度向上に主眼を置いており，候補となった関数対が実行結果の面でも一致するかどうかは十分に検証していない．一方，バイナリの意味比較に関する研究は，関数やループ単位の等価性判定手法を提供するが，TPL検出で得られる大量の候補を対象とした後段検証としての位置づけでは設計されていない．

本研究は，TPL検出等で得られた関数対応候補を入力とし，対象の2バイナリについて対応する入力に対する観測結果を比較して，候補をライブラリ単位へ集約する前にEQまたはNOT\_EQを確認する後段検証器を提案する点で，これらの既存研究と異なる．判定に必要な情報を確認できない場合はUNKNOWNとし，制限時間内に確定できない場合はTIMEOUTとする．
